\documentclass[12pt, a4paper]{article}
\usepackage{../../../myconfig} % Gọi file cấu hình từ ngoài gốc vào

% ==========================================
% BẮT ĐẦU NỘI DUNG TÀI LIỆU
% ==========================================
\begin{document}

% Truyền 3 tham số: {Nhãn cam} {Chữ to chính giữa} {Chữ ở góc phải Header}
\titlebox{Chủ đề: Hàm số}{BTDD hàm doanh thu lợi nhuận chi phí}{Hàm số: Doanh thu lợi nhuận chi phí}

% --- CÂU 1 ---
\questionTrueFalse{0}{1}{Một cơ sở sản xuất dầu ăn đang tính toán các chi phí để vận hành và sản xuất \( x \) lít dầu ăn trong một tháng. Các chi phí bao gồm: 6 triệu đồng chi phí cố định để thuê nhà xưởng, 0,03 triệu đồng chi phí nguyên liệu cần thiết cho mỗi lít dầu ăn, và \( 0,000005x^2 \) triệu đồng chi phí bảo trì máy móc. Biết rằng dầu ăn được bán ra thị trường với giá là 70 nghìn đồng/lít, và giả sử toàn bộ lượng dầu ăn sản xuất được thì đều được bán hết.}{
    \textbf{a)} & Tổng chi phí của cơ sở sản xuất trong một tháng được tính bằng công thức:  
    \( C(x) = 0,000005x^2 + 0,03x + 6 \) (đơn vị: triệu đồng).
    & \textcolor{amberorange}{$\bigcirc$} & \textcolor{amberorange}{$\bigcirc$} \\ \hline
    
    \textbf{b)} & Nếu mỗi tháng cơ sở này sản xuất được 150 lít dầu ăn, khi đó cơ sở sẽ bị lỗ.
    & \textcolor{amberorange}{$\bigcirc$} & \textcolor{amberorange}{$\bigcirc$} \\ \hline
    
    \textbf{c)} & Lợi nhuận lớn nhất mà cơ sở này có thể đạt được trong một tháng lớn hơn 75 triệu đồng.
    & \textcolor{amberorange}{$\bigcirc$} & \textcolor{amberorange}{$\bigcirc$} \\ \hline
    
    \textbf{d)} & Giả sử cơ sở này sản xuất một lượng dầu ăn tiến đến vô cùng, khi đó chi phí trung bình để sản xuất mỗi lít dầu ăn sẽ tiến đến 0.
    & \textcolor{amberorange}{$\bigcirc$} & \textcolor{amberorange}{$\bigcirc$} \\
}

% --- CÂU 2 ---
\questionTrueFalse{0}{2}{Anh Huy là ông chủ của hãng nước hoa Huy's Secret. Anh ước tính rằng nếu anh bán được \( p \) chai nước hoa mỗi tháng (\( p \geq 4 \)), khi đó giá bán của mỗi chai nước hoa (đơn vị: nghìn đồng) được biểu diễn bởi hàm số  
\( B(p) = \dfrac{10000}{p + 75} \)
với tổng chi phí cần thiết để sản xuất \( p \) chai nước hoa được biểu diễn bởi hàm số  
\( C(p) = \dfrac{1}{2000}p^2 + \dfrac{3}{20}p + 400 \)
(đơn vị: nghìn đồng). Xét hàm số \( P(p) \) biểu diễn lợi nhuận mà anh Huy thu được khi bán hết \( p \) chai nước hoa.}{
    \textbf{a)} & Nếu giá bán mỗi chai nước hoa của anh Huy là 100 nghìn đồng/chai, khi đó một tháng anh Huy bán được 25 chai nước hoa.
    & \textcolor{amberorange}{$\bigcirc$} & \textcolor{amberorange}{$\bigcirc$} \\ \hline
    
    \textbf{b)} & Hàm số \( P(p) = C(p) - p \cdot B(p) \).
    & \textcolor{amberorange}{$\bigcirc$} & \textcolor{amberorange}{$\bigcirc$} \\ \hline
    
    \textbf{c)} & anh Huy sẽ bị lỗ nếu bán nhiều hơn 4193 chai nước hoa mỗi tháng.
    & \textcolor{amberorange}{$\bigcirc$} & \textcolor{amberorange}{$\bigcirc$} \\ \hline
    
    \textbf{d)} & Lợi nhuận tối đa mà anh Huy có thể thu về nhờ việc bán nước hoa lớn hơn 8,3 triệu đồng.
    & \textcolor{amberorange}{$\bigcirc$} & \textcolor{amberorange}{$\bigcirc$} \\
}
% --- CÂU 13 ---
\questionShortAnswer{0}{3}{Một khu chung cư có 100 phòng cần cho thuê với mức giá cho mỗi phòng là 5 triệu đồng/tháng. Với mức giá này thì mọi phòng đều có người đến thuê. Để tăng lợi nhuận, chủ khu chung cư dự định sẽ tăng giá thuê thêm 500 nghìn đồng/tháng, tuy nhiên nếu tăng giá như vậy thì sẽ có 8 phòng bị bỏ trống. Hỏi chủ khu chung cư này nên cho thuê bao nhiêu phòng để doanh thu từ việc cho thuê là lớn nhất?}


% --- CÂU 15 ---
\questionShortAnswer{0}{4}{Trong trung tâm thương mại Lotto thuộc thành phố MATHSTER, có một nhà hàng bán buffet nướng lẩu. Khi nhà hàng bán với giá 250 nghìn đồng một suất thì mỗi ngày nhà hàng bán được 100 suất. Nhà hàng dự định có đợt giảm giá để kích cầu nhân dịp đầu năm học mới cho học sinh sinh viên. Theo khảo sát từ thị trường, mỗi lần giảm giá 20 nghìn đồng một suất thì nhà hàng bán thêm được 10 suất. Hỏi nhà hàng cần bán với giá mới là bao nhiêu nghìn đồng một suất để doanh thu trong một ngày là lớn nhất?}
\end{document}